\section{Language-Model-Centered Multimodal Architecture}
\subsection{Foundation and modality-specific encoders}
The full MR-NWM design uses a \emph{pretrained language transformer as the shared
slow backbone}. Continuous observations enter this backbone through trainable
modality adapters rather than being converted into strings. A separate fast
physical subnetwork performs the repeated numeric transitions between backbone
updates. This separation preserves a language-based foundation while assigning
high-frequency prediction to a smaller computational path.

\begin{figure*}[t]
\centering
\begin{tikzpicture}[font=\small,>=Latex,
 box/.style={draw=black!60,rounded corners=2pt,align=center,minimum height=0.9cm},
 arr/.style={->,thick}]
 \node[box,fill=blue!7,text width=3.3cm] (intent) at (0,2.3)
 {Operational preference text};
 \node[box,fill=green!7,text width=3.3cm] (modal) at (0,0)
 {CSI / IQ / KPIs / topology\\Modality encoders\\Age, period, validity};
 \node[box,fill=green!7,text width=3.3cm] (evidence) at (0,-2.3)
 {Observed state and context\\Validated exogenous evidence};
 \node[box,fill=violet!9,text width=3.4cm,minimum height=3.1cm] (lm) at (5,0)
 {\textbf{Shared pretrained\\language transformer}\\[0.15cm]
 Physical view:\\observations + evidence\\[0.1cm]
 Semantic view:\\context + intent};
 \node[box,fill=orange!10,text width=3.5cm] (packet) at (10,1.5)
 {Context projection\\Cached FiLM / low-rank packet};
 \node[box,fill=orange!10,text width=3.5cm] (fast) at (10,-0.2)
 {Fast recurrent physics\\State, action, memory\\$K$ steps per slow refresh};
 \node[box,fill=blue!7,text width=3.4cm] (weights) at (5,-2.7)
 {Structured objective weights};
 \node[box,fill=orange!10,text width=3.5cm] (states) at (10,-2.2)
 {Continuous state trajectory\\Summary to next slow update};
 \draw[arr] (intent.east) -- (lm.north west);
 \draw[arr] (modal.east) -- (lm.west);
 \draw[arr] (evidence.east) -- (lm.south west);
 \draw[arr] (lm.north east) -- (packet.west);
 \draw[arr] (packet) -- (fast);
 \draw[arr] (fast) -- (states);
 \draw[arr] (lm.south) -- (weights.north);
\end{tikzpicture}
\caption{Language-centered architecture implemented by the foundation path. The
language transformer is the shared slow backbone; two separate views enforce the
preference/evidence boundary. Modality adapters and wireless heads require task
training. Fast predictions use the cached packet without calling the language model.}
\label{fig:language-foundation}
\end{figure*}

Let $x_T^{(m)}$ be observations of modality $m$, $v_T^{(m)}$ their validity mask,
$\delta_T^{(m)}$ their observation ages, and $\rho_m$ the acquisition period.
A modality encoder creates continuous tokens
\begin{equation}
 e_T^{(m)}=E_m(x_T^{(m)})+e_{\rm type}^{(m)}
       +E_{\rm time}[\log(1+\delta_T^{(m)}),\log(1+\rho_m)].
\end{equation}
The implementation supports variable token counts, explicitly masked missing
observations, and an input-token budget. A modality can be absent entirely.
Negative ages, non-finite valid measurements, and malformed masks are rejected.
These checks make the temporal information explicit; they are not a distributed
clock-synchronization or missing-data recovery system.

\begin{table*}[t]
\centering\small
\caption{Multimodal inputs in the language-foundation integration. Front ends are
trainable adapters; their radio representations require task-specific data and alignment.}
\begin{tabular}{p{0.16\textwidth}p{0.35\textwidth}p{0.40\textwidth}}
\toprule
Input & Representation entering the model & Implemented boundary and limitations\\
\midrule
Intent, protocol/evidence text & Language tokenizer IDs and masks & Preferences and exogenous evidence use separate typed fields.\\
State, traffic, KPIs & Continuous vectors or sequences & Numeric projections preserve values; timestamps and acquisition periods are supplied by the caller.\\
CSI and IQ & Continuous tokens with real/imaginary feature pairs & Linear/MLP adapters; complex preprocessing and RF calibration are external.\\
Topology & Node features and optional adjacency & One normalized graph message-passing step before language-backbone fusion.\\
Visual/map features & Continuous patch feature vectors & Configurable numeric patch adapter; raw-image patch extraction is external.\\
\bottomrule
\end{tabular}
\label{tab:modalities}
\end{table*}

For graph inputs, an optional normalized adjacency operator mixes neighboring
node embeddings before fusion. For CSI/IQ, the caller supplies explicit real
and imaginary components with known dimensions. Image or map patches can use
the same registered-adapter interface, but this release does not include a
pretrained vision encoder or raw-camera pipeline. Adapter presence alone does
not establish useful cross-modal alignment. That alignment must be learned and
tested on synchronized observation/action trajectories.

\subsection{Two views of one language backbone}
Physical predictions must be able to use information about the external world
without confusing it with the operator's preferences. The same pretrained
transformer $F_\theta$ therefore processes two separate input views. Let
\begin{equation}
 Z_T^p=[E_s(s_{T,0});E_c(c_T);E_H(H_{T-1});
                  \{e_T^{(m)}\}_m;E_K(K_T)],
\end{equation}
where $K_T$ contains externally supplied, validated evidence, such as an actual
configuration change. A learned query $q_p$ reads the physical view, while
intent tokens are appended only to the semantic view:
\begin{align}
 h_T^p&=P_pF_\theta([Z_T^p;q_p])_{\rm last},\\
 h_T^I&=P_IF_\theta([Z_T^p;E_I(I_T)])_{\rm last\ valid}.
\end{align}
Both views use the \emph{same language-model parameters}, with independent
forward passes and masks. Preferences cannot enter the physical forward pass
through attention or hidden-state reuse. By contrast, changes to observed
modalities or exogenous evidence may legitimately change predicted dynamics.
The caller is responsible for distinguishing a requested change from a change
that has actually been applied or externally observed.

The projected physical context passes through the recurrent physics expert and
then generates the fast-tier FiLM/low-rank packet. The physics expert can share
its parameters with the fast recurrence. This is a coupling around the language
trunk; the code does not claim to have converted the pretrained transformer's
internal feed-forward layers into wireless SSM experts. The objective head
uses $h_T^I$. Continuous state and structured objective heads are implemented;
general natural-language explanations and arbitrary protocol action decoders
remain additional output heads to develop and evaluate.

The \texttt{LanguageBasedWorldModel.from\_pretrained} entry point loads an actual
Hugging Face transformer checkpoint. A concrete small example is
\texttt{Qwen/Qwen2.5-0.5B} \cite{qwen2024}. Pretrained weights can be frozen while
gradients still propagate \emph{through} the transformer into numeric adapters.
Disabling autograd around the frozen trunk would break that training path and
is deliberately avoided. The numeric adapters, context projections, modulation
heads, and dynamics predictor start untrained; pretrained text weights do not
automatically provide a calibrated wireless world model.

\subsection{Multiple observation rates and execution rates}
The architecture distinguishes three quantities: each modality's acquisition
period $\rho_m$, the slow backbone's refresh interval $\Delta T$, and the fast
predictor's step interval $\Delta t$, with $\Delta T=K\Delta t$ in the basic
training graph. Observation tokens retain their own ages and periods even when
combined at one slow refresh. A token budget requires explicit pooling or
resampling before excessively long streams enter the language backbone.

Only two learned execution rates are implemented: one language-backbone update
per macro window and $K$ lightweight recurrent steps. Operator-level knowledge
retrieval, long-horizon task decomposition, and model lifecycle management can
run at a still slower external rate, but are not a third implemented neural
controller in this release. The fast API does not call the language backbone;
the contract tests count backbone calls to enforce this property.

Placement in Non-RT management, a Near-RT application, or an in-node predictor
depends on measured end-to-end budgets. No general language-model size or Python
module name establishes that placement. In particular, the language-foundation
version's latency must be measured separately from the lightweight prototype's
timing table. RF/CSI sampling, tensor fusion, model inference, network transport,
and actuator execution contribute different parts of the overall budget.

\subsection{What the two code paths establish}
The \texttt{foundation.py} path implements the intended language-centered
architecture and registered multimodal token adapters. Its tests cover causal
preference separation, graph/missing-data handling, input ages, frozen-trunk
gradient flow, and cached fast execution. The optional integration script records
the checkpoint identity and resolved revision. We executed this check with
Qwen2.5-0.5B on CPU, loading 494,032,768 pretrained backbone parameters and
connecting CSI, IQ, and topology feature tokens. The five-step output had shape
$1\times5\times6$; numeric-adapter gradients were nonzero with frozen backbone
weights, and preference changes left physical predictions unchanged. The record
is stored in \texttt{pretrained\_integration.json}. This establishes
software/gradient compatibility rather than a trained RAN model or a general
language-understanding score.

The following lightweight variant replaces the expensive language trunk with
small train-from-scratch attention/MLP components. Its controlled experiments
isolate recurrent dynamics, sharing, and modulation at low cost. The numerical
results later in this paper belong \emph{only} to that lightweight variant.
They are not evaluations of Qwen or evidence that language pretraining improves
wireless prediction.
